% !TEX program = pdflatex
% =============================================================================
%  RaicesMX · Proyectos I+D · Equipo 10 · BUAP
%  ENTREGA 1 EN FORMATO IEEE (una columna)
%  Documento del modelo E-R implementado y de los 5 casos de uso
%  implementados en la primera versión del sistema.
%
%  ------------------------------------------------------------------
%  CÓMO USARLO EN OVERLEAF
%  ------------------------------------------------------------------
%  1. Proyecto nuevo en Overleaf, compilador pdfLaTeX.
%  2. Pega todo este contenido en main.tex.
%  3. Sube las imágenes en una carpeta "figuras/" con estos nombres:
%        logo-buap.png (o escudo_positivo.png)
%        modelo-er.png        (diagrama E-R)
%        cu-01-registro.png   cu-04-publicar.png   cu-06-verificar.png
%        cu-07-catalogo.png   cu-08-compra.png
%     Si falta alguna, el documento dibuja un marcador con el nombre y
%     compila igual.
%  4. Compila. Listo.
%
%  NOTAS
%  - Hoja tamaño carta (letterpaper), estándar de IEEE.
%  - Las referencias van numeradas estilo IEEE, pero conservan los datos
%    completos en formato APA del equipo.
%  - Portada con la paleta de RaicesMX: alabastro #F2F0EB, verde milpa
%    #2F6B3B, verde oscuro #24512E, tinta #2B211C, verde claro #B9DCAE.
% =============================================================================

\documentclass[conference,letterpaper,onecolumn]{IEEEtran}

\usepackage[utf8]{inputenc}
\usepackage[T1]{fontenc}
\usepackage[spanish,es-nodecimaldot,es-noquoting,es-tabla]{babel}

\usepackage{graphicx}
\usepackage[table]{xcolor}
\usepackage{booktabs}
\usepackage{array}
\usepackage{amsmath}
\usepackage{microtype}
\usepackage{multirow}
\usepackage{float}   % [H] deja cada figura exactamente donde se escribe
\usepackage[hyphens]{url}
\usepackage{url}
\usepackage{hyperref}

\newcolumntype{L}[1]{>{\raggedright\arraybackslash}p{#1}}

\definecolor{VerdeMilpa}{HTML}{2F6B3B}
\definecolor{VerdeOscuro}{HTML}{24512E}
\definecolor{Alabastro}{HTML}{F2F0EB}
\definecolor{Tinta}{HTML}{2B211C}
\definecolor{VerdeClaro}{HTML}{B9DCAE}
\definecolor{GrisMedio}{HTML}{8A817A}

\hypersetup{
  colorlinks=true,
  linkcolor=VerdeMilpa,
  urlcolor=VerdeOscuro,
  citecolor=VerdeMilpa,
  pdftitle={RaicesMX: modelo E-R implementado y cinco casos de uso},
  pdfauthor={Equipo 10 - Proyectos I+D - BUAP}
}

% =============================================================================
%  DATOS DEL PROYECTO  ·  se editan únicamente aquí
% =============================================================================
\newcommand{\Proyecto}{RaicesMX}
\newcommand{\Materia}{Proyectos I+D}
\newcommand{\Equipo}{Equipo 10}
\newcommand{\Titular}{Mtra. S\'anchez Rom\'an Guillermina}
\newcommand{\Periodo}{Oto\~no 2026}
\newcommand{\Institucion}{Benem\'erita Universidad Aut\'onoma de Puebla}
\newcommand{\Facultad}{Facultad de Ciencias de la Computaci\'on}
% =============================================================================
%  FIN DE LOS DATOS DEL PROYECTO
% =============================================================================

% ---------------- Figura con marcador si falta el archivo ----------------
% Todas las figuras y tablas van con [H] para quedar exactamente donde se
% escriben (formato de una columna).
\newcommand{\figura}[4]{%
  \begin{figure}[H]
  \centering
  \IfFileExists{#1}{%
    \includegraphics[width=#3]{#1}%
  }{%
    \begingroup\setlength{\fboxsep}{6pt}%
    \fcolorbox{GrisMedio}{Alabastro}{%
      \parbox[c][3.0cm][c]{\dimexpr#3-2\fboxsep-2\fboxrule\relax}{%
        \centering\color{GrisMedio}\footnotesize
        \textbf{[ IMAGEN ]}\\[3pt]\texttt{\detokenize{#1}}}%
    }%
    \endgroup
  }
  \caption{#2}\label{#4}
  \end{figure}%
}

\begin{document}

% ============================== PORTADA ==================================
\pagecolor{Alabastro}
\thispagestyle{empty}

\noindent{\color{VerdeMilpa}\rule{\textwidth}{3pt}}\\[-6pt]
\noindent{\color{VerdeClaro}\rule{\textwidth}{1.2pt}}

\vspace{0.22cm}

\begin{center}
\IfFileExists{figuras/logo-buap.png}{%
  \includegraphics[width=2.7cm]{figuras/logo-buap.png}%
}{%
  \IfFileExists{logo-buap.png}{%
    \includegraphics[width=2.7cm]{logo-buap.png}%
  }{%
    \IfFileExists{escudo_positivo.png}{%
      \includegraphics[width=2.7cm]{escudo_positivo.png}%
    }{%
      \begingroup\setlength{\fboxsep}{6pt}%
      \fcolorbox{GrisMedio}{white}{\parbox[c][2.2cm][c]{3.4cm}{\centering\color{GrisMedio}\scriptsize LOGO\\BUAP}}%
      \endgroup
    }%
  }%
}

\vspace{0.22cm}
{\color{GrisMedio}\small \Institucion\\\Facultad\par}
{\color{GrisMedio}\scriptsize \Materia\ \textbullet\ \Equipo\ \textbullet\ \Periodo\par}
\end{center}

\vspace{0.25cm}

\begin{center}
{\color{VerdeMilpa}\Huge\bfseries \Proyecto\par}
\vspace{0.15cm}
{\color{VerdeOscuro}\large
Comercialización directa para productores originarios mexicanos
con verificación criptográfica del origen de la pieza\par}
\vspace{0.25cm}
{\color{GrisMedio}\normalsize
Entrega 1: modelo E-R de la base de datos (implementación) y cinco casos
de uso implementados en la primera versión del sistema\par}
\end{center}

\vspace{0.22cm}

\begin{center}
\begingroup
\setlength{\fboxsep}{10pt}
\fcolorbox{VerdeMilpa!45}{VerdeMilpa!8}{%
\begin{minipage}{0.90\textwidth}
\renewcommand{\arraystretch}{1.25}
\begin{tabular}{@{}>{\bfseries\color{VerdeMilpa}}l@{\hspace{1.4em}}l@{}}
Proyecto & \Proyecto, plataforma web progresiva (PWA) de venta directa \\
Asignatura & \Materia\ (proyecto de software) \\
Titular & \Titular \\
Equipo & \Equipo\ (cinco integrantes) \\
Periodo & \Periodo \\
Alcance & Modelo E-R implementado y 5 casos de uso en la versión~1.0 \\
Stack & Next.js 16, React 19, TypeScript, Tailwind CSS, PostgreSQL 15 (Supabase) \\
\end{tabular}
\end{minipage}}
\endgroup
\end{center}

\vspace{0.22cm}

% =============================================================================
%  PORTADA · INTEGRANTES (una persona por línea)
% =============================================================================
\begin{center}
\begingroup
\setlength{\fboxsep}{11pt}
\fcolorbox{VerdeMilpa!45}{VerdeMilpa!8}{%
\begin{minipage}{0.82\textwidth}
\centering
{\color{VerdeMilpa}\bfseries\small INTEGRANTES}\\[9pt]
\renewcommand{\arraystretch}{1.55}
\begin{tabular}{@{}>{\color{VerdeMilpa}\bfseries}l@{\hspace{0.9em}}L{0.58\linewidth}r@{}}
1 & Ángel Armando Carreño González & 202108258 \\
2 & Miguel Ángel Hernández Sánchez & 202123950 \\
3 & Alan Espinoza Flores & 202260523 \\
4 & Luis Ángel Zacarías Ávila & 202166331 \\
5 & Adrián Torres Rodríguez & 202163210 \\
\end{tabular}
\end{minipage}}
\endgroup
\end{center}

\vspace{0.3cm}

\begin{center}
\begin{minipage}{0.94\textwidth}
\centering\footnotesize
\textbf{Contenido:} I. Introducción --- II. Modelo E-R de la base de datos
(implementación) --- III. Cinco casos de uso implementados en la versión~1.0 ---
IV. Seguridad de la base --- V. Pruebas --- VI. Conclusiones --- Referencias.
\end{minipage}
\end{center}

\vspace{0.6cm}

\begin{center}
{\color{GrisMedio}\footnotesize
\Institucion\ \textbullet\ \Facultad\ \textbullet\ \Periodo\par}
{\color{VerdeMilpa}\small Documento de entrega 1 \textbullet\ Formato IEEE (una columna)\par}
\end{center}

\noindent{\color{VerdeClaro}\rule{\textwidth}{1.2pt}}\\[-6pt]
\noindent{\color{VerdeMilpa}\rule{\textwidth}{3pt}}

\nopagecolor

% ============================ ARTÍCULO IEEE ==============================
\title{\Proyecto: modelo entidad-relación implementado y cinco casos de uso
en la primera versión del sistema}

\author{%
\IEEEauthorblockN{Ángel A. Carreño González\and
Miguel Á. Hernández Sánchez\and
Alan Espinoza Flores}
\IEEEauthorblockN{Luis Á. Zacarías Ávila\and
Adrián Torres Rodríguez}}

\maketitle

\begin{abstract}
\noindent
Este documento entrega los dos productos solicitados en la primera entrega de
\Proyecto: el \textbf{modelo entidad-relación de la base de datos, implementado} en
PostgreSQL, y la \textbf{implementación de cinco casos de uso} en la versión~1.0 del
sistema. El modelo consta de 17 tablas, 10 tipos enumerados, 2 vistas, 1 función
transaccional y 8 disparadores, y cubre las 24 clases del diagrama UML del proyecto;
las reglas de negocio que antes residirían en la aplicación quedaron como
restricciones \texttt{CHECK}, disparadores y vistas de la propia base. La versión~1.0
del sistema es un monolito modular sobre Next.js 16 con una API REST única, y sostiene
los cinco casos de uso ---registro y verificación de productores, publicación de
piezas con ficha de origen y sello QR firmado, verificación pública del sello,
catálogo con filtros y mapa de procedencia, y compra con reparto 90/10--- operados
contra la base real, no en memoria. Se documenta la trazabilidad entre el modelo UML
y el esquema, las decisiones de seguridad a nivel de fila y las evidencias de cada
prueba ejecutada.
\end{abstract}

\begin{IEEEkeywords}
modelo entidad-relación, PostgreSQL, casos de uso, comercio electrónico, artesanías,
verificación criptográfica, seguridad a nivel de fila.
\end{IEEEkeywords}

% ============================== I. INTRODUCCIÓN ==========================
\section{Introducción}

\Proyecto\ es una plataforma web progresiva para que los productores originarios
mexicanos vendan sus piezas directo, sin intermediarios. Cada pieza publica nueve
datos de origen obligatorios y un sello QR firmado que cualquier persona puede
comprobar sin crear una cuenta; en cada venta el 90\,\% es para el taller y el
10\,\% para la plataforma, y ese desglose se muestra al comprador antes de pagar.

La primera entrega del proyecto comprende exactamente dos productos, y este documento
los presenta en ese orden: \emph{(i)} el modelo entidad-relación de la base de datos
\textbf{implementado} en el motor, con su trazabilidad contra el diagrama de clases
UML y con las reglas de integridad residentes en la base; y \emph{(ii)} la
implementación de \textbf{cinco casos de uso} en la versión~1.0 del sistema, cada uno
con su pantalla, su API, su flujo y la evidencia de su ejecución contra la base real.

El sistema se construyó con Next.js 16, React 19, TypeScript y Tailwind CSS, con
PostgreSQL 15 administrado en Supabase y una API REST única bajo \texttt{/api/v1};
el desarrollo se organizó con Scrum en sprints de dos semanas \cite{propuesta}. La
versión~1.0 no se despliega: corre en local contra la base en la nube.

% ===================== II. MODELO E-R (IMPLEMENTACIÓN) ====================
\section{Modelo E-R de la base de datos (implementación)}

El modelo se implementó en PostgreSQL mediante tres migraciones versionadas
(\texttt{001} esquema inicial, \texttt{002} reglas adicionales y \texttt{003}
notificaciones y reportes) y se aplicó con un aprovisionamiento idempotente
(\texttt{npm run db:setup}) que además carga los datos semilla. La Figura
\ref{fig:er} muestra el modelo tal como quedó en el motor.

\figura{figuras/modelo-er.png}{Modelo entidad-relación implementado en PostgreSQL: 17 tablas, 10 enumeraciones, 2 vistas, 1 función transaccional y 8 disparadores. La versión en alta resolución (3200$\times$2788) está en \texttt{docs/modelo-er.png} del repositorio.}{0.95\textwidth}{fig:er}

El esquema se organiza en cuatro dominios y las subclases del UML se materializan
como especializaciones por rol sobre \texttt{usuarios}, con las tablas
\nolinkurl{perfiles_productor} y \nolinkurl{perfiles_comprador} para los atributos
propios de cada perfil. La Tabla \ref{tab:dominios} resume la composición del modelo.

\begin{table}[H]
\caption{Composición del modelo implementado.}
\label{tab:dominios}
\centering
\renewcommand{\arraystretch}{1.08}
\begin{tabular}{@{}L{0.30\columnwidth}L{0.62\columnwidth}@{}}
\toprule
\textbf{Dominio} & \textbf{Tablas} \\
\midrule
Usuarios y verificación & \texttt{usuarios}, \nolinkurl{perfiles_productor}, \nolinkurl{verificaciones_productor}, \nolinkurl{perfiles_comprador} \\
Catálogo y procedencia & \texttt{productos}, \nolinkurl{fichas_origen}, \nolinkurl{sellos_originalidad} \\
Pedidos, pagos y acuerdos & \texttt{pedidos}, \nolinkurl{detalles_pedido}, \nolinkurl{contratos_b2b}, \texttt{pagos}, \texttt{liquidaciones} \\
Seguimiento y control & \nolinkurl{seguimiento_pedidos}, \texttt{calificaciones}, \texttt{notificaciones}, \texttt{reportes}, \nolinkurl{bitacora_auditoria} \\
\midrule
\multicolumn{2}{@{}L{0.94\columnwidth}@{}}{\texttt{usuarios}, \texttt{pedidos} y \texttt{productos} concentran las 10 enumeraciones; \nolinkurl{vista_catalogo} y \nolinkurl{vista_mapa_procedencia} son vistas de lectura} \\
\bottomrule
\end{tabular}
\end{table}

La Tabla \ref{tab:traza} muestra la trazabilidad completa contra el diagrama de
clases UML del proyecto \cite{uml}: las 24 clases tienen correspondencia directa, ya
sea como tabla, como columna calculada o como especialización por rol.

\begin{table}[H]
\caption{Trazabilidad clase UML $\rightarrow$ implementación (24 clases).}
\label{tab:traza}
\centering
\footnotesize
\renewcommand{\arraystretch}{1.05}
\begin{tabular}{@{}L{0.19\textwidth}L{0.25\textwidth}L{0.19\textwidth}L{0.22\textwidth}@{}}
\toprule
\textbf{Clase UML} & \textbf{Implementación} & \textbf{Tipo} & \textbf{Notas} \\
\midrule
\texttt{Usuario} (abstracta) & \texttt{usuarios} + enum \texttt{rol\_usuario} & Tabla & Identidad y credenciales \\
\texttt{Productor} & \nolinkurl{perfiles_productor} & Tabla & PK compartida con \texttt{usuarios} \\
\texttt{Comprador} (abstracta) & \nolinkurl{perfiles_comprador} + rol & Tabla & Especialización por rol \\
\texttt{CompradorB2C} & \texttt{usuarios.rol = comprador\_b2c} & Especialización & Comprador final \\
\texttt{CompradorB2B} & \texttt{usuarios.rol = comprador\_b2b} & Especialización & Requiere empresa \\
\texttt{Promotor} & \texttt{usuarios.rol = promotor} & Especialización & Verifica productores \\
\texttt{Administrador} & \texttt{usuarios.rol = admin} & Especialización & Dos hojas del UML \\
\texttt{Producto} & \texttt{productos} & Tabla & Precio, stock y categoría \\
\texttt{FichaOrigen} & \nolinkurl{fichas_origen} & Tabla & 1:1 con producto, 9 campos \\
\texttt{SelloOriginalidad} & \nolinkurl{sellos_originalidad} & Tabla & Folio + HMAC-SHA256 \\
\texttt{VerificacionProductor} & \nolinkurl{verificaciones_productor} & Tabla & Registro de revisión \\
\texttt{Pedido} (abstracta) & \texttt{pedidos} + enum \texttt{tipo\_pedido} & Tabla & Reparte el dinero \\
\texttt{PedidoB2C} & \texttt{pedidos.tipo = B2C} & Especialización & Comisión 10\,\% \\
\texttt{PedidoB2B} & \texttt{pedidos.tipo = B2B} + \nolinkurl{contratos_b2b} & Especialización & Comisión 5\,\%, pago 30/70 \\
\texttt{DetallePedido} & \nolinkurl{detalles_pedido} & Tabla & Línea por producto \\
\texttt{Pago} & \texttt{pagos} & Tabla & Tarjeta o SPEI (simulado) \\
\texttt{Comision} & \nolinkurl{pedidos.monto_comision} & Columnas calculadas & Se deriva del pedido \\
\texttt{Liquidacion} & \texttt{liquidaciones} & Tabla & Neto del productor \\
\texttt{ContratoB2B} & \nolinkurl{contratos_b2b} & Tabla & Anticipo 30 / saldo 70 \\
\texttt{SeguimientoPedido} & \nolinkurl{seguimiento_pedidos} & Tabla & Historial automático \\
\texttt{Calificacion} & \texttt{calificaciones} & Tabla & 1--5, una por pedido \\
\texttt{Notificacion} & \texttt{notificaciones} & Tabla & Migración 003 \\
\texttt{Reporte} & \texttt{reportes} & Tabla & Migración 003 \\
\texttt{BitacoraAuditoria} & \nolinkurl{bitacora_auditoria} & Tabla & Solo \texttt{INSERT} \\
\bottomrule
\end{tabular}
\end{table}

Una decisión de diseño relevante es que la clase \texttt{Comision} no se materializa
como entidad independiente porque siempre se deriva del pedido: se conservan el monto
y el porcentaje en \texttt{pedidos} y el desglose completo en \texttt{liquidaciones},
con una restricción que obliga a que productor más comisión sumen exactamente lo
vendido. La Tabla \ref{tab:reglas} resume las reglas que la base defiende por sí
sola.

\begin{table}[H]
\caption{Reglas de integridad implementadas en el motor.}
\label{tab:reglas}
\centering
\renewcommand{\arraystretch}{1.05}
\begin{tabular}{@{}L{0.50\columnwidth}L{0.42\columnwidth}@{}}
\toprule
\textbf{Regla} & \textbf{Dónde vive} \\
\midrule
La ficha exige los 9 campos de origen & \texttt{CHECK} por columna en \nolinkurl{fichas_origen} \\
Solo un productor aprobado publica & Disparador \nolinkurl{trg_productos_productor_aprobado} \\
Productor + comisión = subtotal $-$ descuento & \texttt{CHECK ck\_pedido\_reparto} \\
B2B solo de 5 a 50 piezas & \texttt{CHECK ck\_pedido\_lote} \\
Compra concurrente segura & Función \texttt{crear\_pedidos\_b2c} con \texttt{FOR UPDATE} \\
La bitácora es inmutable & Disparadores \texttt{BEFORE UPDATE/DELETE/TRUNCATE} \\
Búsqueda sin acentos & Función \texttt{sin\_acentos} + vista \texttt{vista\_catalogo} \\
Historial y promedio automáticos & Disparadores de seguimiento y calificaciones \\
\bottomrule
\end{tabular}
\end{table}

% ================ III. CINCO CASOS DE USO IMPLEMENTADOS =================
\section{Cinco casos de uso implementados en la versión~1.0}

Los cinco casos de uso seleccionados cubren el flujo completo, desde que el artesano se
registra hasta que alguien le compra: \textbf{CU-01} registro y verificación del
productor, \textbf{CU-04} publicación de pieza con ficha de origen y sello QR,
\textbf{CU-06} verificación del sello sin cuenta, \textbf{CU-07} catálogo con filtros
y mapa de procedencia, y \textbf{CU-08} carrito y compra con reparto 90/10. La
Tabla \ref{tab:cu} resume su implementación; después se documenta cada uno.

\begin{table}[H]
\caption{Los cinco casos de uso de la versión~1.0.}
\label{tab:cu}
\centering
\renewcommand{\arraystretch}{1.05}
\begin{tabular}{@{}L{0.10\columnwidth}L{0.38\columnwidth}L{0.42\columnwidth}@{}}
\toprule
\textbf{CU} & \textbf{Pantalla} & \textbf{Evidencia en la base} \\
\midrule
CU-01 & \texttt{/registro-productor}, \texttt{/promotor} & Productor \texttt{9da68be6-}\dots\ creado pendiente y aprobado \\
CU-04 & \texttt{/productor/publicar} & Producto \texttt{1fb66f12-}\dots, folio \nolinkurl{RMX-PUE-2026-1EA2E8} \\
CU-06 & \texttt{/verificar/[folio]} & \texttt{verificado: true} en \nolinkurl{RMX-PUE-2026-3F9A1C} \\
CU-07 & \texttt{/catalogo} & Puebla 4 productores / 7 piezas; Oaxaca 1 / 1 \\
CU-08 & \texttt{/comprador/carrito} & \texttt{PED-2026-C1C0A80D}: \$350 + \$35 + \$150 = \$500 \\
\bottomrule
\end{tabular}
\end{table}

\subsection{CU-01: registro y verificación de productores}

\emph{Actores:} productor y promotor.\\
\emph{Pantallas:} \nolinkurl{/registro-productor} y \nolinkurl{/promotor}.\\
\emph{API:} \nolinkurl{POST /api/v1/productores} y \nolinkurl{PATCH /api/v1/productores/:id/verificacion}.

\emph{Flujo:} (1) el productor captura sus datos generales y la ficha artesanal con
CURP o clave de elector; (2) el servidor valida los campos obligatorios; (3) se
insertan \texttt{usuarios} y \nolinkurl{perfiles_productor} con estado
\texttt{pendiente}; (4) el promotor aprueba o rechaza desde el panel, y el sistema
registra la verificación y la bitácora; (5) solo un productor \texttt{aprobado} puede
publicar, por disparador de la base.

\emph{Evidencia:} el productor \nolinkurl{9da68be6-cce2-4a44-9a72-0c8509a75a80} fue
creado en estado \texttt{pendiente} y aprobado después; el conteo de aprobados pasó
de 5 a 6 y el cambio quedó en \nolinkurl{verificaciones_productor} y en la bitácora.

La Figura~\ref{fig:cu01} muestra el formulario de registro con la ficha artesanal y
el panel donde el promotor la revisa.

\figura{figuras/cu-01-registro.png}{Formulario de registro del productor y panel del promotor para su verificación.}{0.95\columnwidth}{fig:cu01}

\subsection{CU-04: publicación de pieza con ficha de origen y sello QR}

\emph{Actor:} productor aprobado.\\
\emph{Pantalla:} \nolinkurl{/productor/publicar}.\\
\emph{API:} \nolinkurl{POST /api/v1/productos} (\texttt{multipart/form-data}).

\emph{Flujo:} (1) el productor captura los nueve campos obligatorios de la ficha, el
precio y el stock; (2) el servidor valida y genera el folio
\texttt{RMX-PUE-2026-XXXXXX}; (3) se calcula la firma HMAC-SHA256 sobre la ficha;
(4) se insertan producto, ficha y sello, y se sube la fotografía al bucket
\texttt{fotos-piezas}; (5) se devuelve el QR y la URL pública de verificación.

\emph{Evidencia:} producto \nolinkurl{1fb66f12-d382-4c05-ba3c-d56fd13e120f}, ficha
\nolinkurl{6d1294f4-c38e-4d06-b193-89932e941e81}, folio \textbf{RMX-PUE-2026-1EA2E8} y
firma \texttt{9cc9d8d4}\dots\texttt{61da}.

La Figura~\ref{fig:cu04} muestra el formulario de publicación y el sello QR que se
genera al guardar la pieza.

\figura{figuras/cu-04-publicar.png}{Formulario de publicación y sello QR generado con el folio y la firma de la pieza.}{0.95\columnwidth}{fig:cu04}

\subsection{CU-06: verificación del origen de una pieza por QR}

\emph{Actores:} comprador o público sin cuenta.\\
\emph{Pantalla:} \nolinkurl{/verificar/[folio]}.\\
\emph{API:} \nolinkurl{GET /api/v1/sello-qr/:folio?sig=} (pública).

\emph{Flujo:} (1) el comprador escanea el QR o abre el enlace; (2) el servidor
recalcula el HMAC con el payload de la ficha; (3) compara con la firma guardada y
muestra la ficha completa o la alerta de falsificación.

\emph{Evidencia:} los folios \nolinkurl{RMX-PUE-2026-3F9A1C} (semilla) y
\nolinkurl{RMX-PUE-2026-1EA2E8} (emitido en la prueba) responden
\texttt{verificado: true}. Si alguien altera la pieza, la firma deja de coincidir.

La Figura~\ref{fig:cu06} muestra la pantalla pública que cualquiera puede abrir
con el QR, sin crear cuenta.

\figura{figuras/cu-06-verificar.png}{Pantalla pública de verificación con la etiqueta ``Origen verificado'' y la ficha completa.}{0.95\columnwidth}{fig:cu06}

\subsection{CU-07: catálogo con filtros y mapa de procedencia}

\emph{Actor:} comprador.\\
\emph{Pantalla:} \nolinkurl{/catalogo}.\\
\emph{API:} \nolinkurl{GET /api/v1/productos} con filtros de categoría, estado, pueblo,
técnica y precio, y \nolinkurl{GET /api/v1/mapa}.

\emph{Flujo:} (1) el catálogo lee de la vista \texttt{vista\_catalogo}, que expone
solo piezas activas con existencias y de productores aprobados; (2) los filtros se
aplican en SQL ignorando acentos; (3) el mapa agrupa por estado en
\texttt{vista\_mapa\_procedencia} y cada casilla filtra el catálogo.

\emph{Evidencia:} el mapa reporta Puebla con 4 productores y 7 piezas, y Oaxaca con 1
productor y 1 pieza; el catálogo muestra 8 piezas antes de la prueba.

La Figura~\ref{fig:cu07} muestra el catálogo con sus filtros y el mapa nacional de
procedencia.

\figura{figuras/cu-07-catalogo.png}{Catálogo con filtros y mapa nacional de procedencia.}{0.95\columnwidth}{fig:cu07}

\subsection{CU-08: compra con desglose transparente 90/10}

\emph{Actor:} comprador B2C.\\
\emph{Pantalla:} \nolinkurl{/comprador/carrito}.\\
\emph{API:} \nolinkurl{POST /api/v1/pedidos/cotizacion} y \nolinkurl{POST /api/v1/pedidos}.

\emph{Flujo:} (1) la cotización usa siempre los precios de la base, nunca los que
envía el cliente; (2) el carrito agrupa por productor y muestra el reparto; (3) al
confirmar se ejecuta la función transaccional \texttt{crear\_pedidos\_b2c}, que
descuenta existencias y crea pago, liquidación, seguimiento y bitácora en una sola
transacción.

\emph{Evidencia:} la Tabla \ref{tab:cu08} muestra el desglose de la compra de prueba
registrada en la base, y la Figura~\ref{fig:cu08} el carrito con el reparto antes de
pagar.

\begin{table}[H]
\caption{Desglose de la compra de prueba registrada en la base.}
\label{tab:cu08}
\centering
\renewcommand{\arraystretch}{1.05}
\begin{tabular}{@{}L{0.58\columnwidth}r@{}}
\toprule
\textbf{Concepto} & \textbf{Valor} \\
\midrule
Folio del pedido & \texttt{PED-2026-C1C0A80D} \\
Subtotal & \$350.00 \\
Comisión de la plataforma (10\,\%) & \$35.00 \\
Neto del productor (90\,\%) & \$315.00 \\
Envío & \$150.00 \\
\textbf{Total cobrado} & \textbf{\$500.00} \\
Existencias del tenate & 10 $\rightarrow$ 9 \\
\bottomrule
\end{tabular}
\end{table}

\figura{figuras/cu-08-compra.png}{Carrito con el desglose transparente del reparto entre plataforma y taller.}{0.95\columnwidth}{fig:cu08}

% ============================ IV. SEGURIDAD ==============================
\section{Seguridad de la base de datos}

La Tabla \ref{tab:seguridad} resume las decisiones de seguridad de la versión~1.0.

\begin{table}[H]
\caption{Decisiones de seguridad.}
\label{tab:seguridad}
\centering
\renewcommand{\arraystretch}{1.05}
\begin{tabular}{@{}L{0.48\columnwidth}L{0.44\columnwidth}@{}}
\toprule
\textbf{Decisión} & \textbf{Por qué} \\
\midrule
RLS activo en las 17 tablas, sin políticas para \texttt{anon} & La llave pública no lee ni escribe aunque se filtre \\
La llave secreta solo en el servidor & El navegador nunca habla directo con PostgreSQL \\
\texttt{revoke all} sobre vistas y tablas nuevas & Cierra el acceso por PostgREST \\
Bitácora inmutable por disparador & Auditoría que la aplicación no puede alterar \\
Folio y firma generados en el servidor & El cliente no puede falsificar un sello \\
Precios tomados de la base & El comprador no manipula el importe \\
\bottomrule
\end{tabular}
\end{table}

La firma del folio es HMAC-SHA256 \cite{rfc2104}, que detecta cualquier alteración
del contenido firmado sin necesidad de certificados; el control de acceso sigue el
principio de mínimo privilegio \cite{owasp} y la seguridad a nivel de fila de
PostgreSQL \cite{rowlevel}. La bitácora acepta únicamente \texttt{INSERT} y sus
disparadores bloquean \texttt{UPDATE}, \texttt{DELETE} y \texttt{TRUNCATE}.

% ============================== V. PRUEBAS ===============================
\section{Pruebas}

La verificación combinó pruebas unitarias sobre la lógica más delicada, resumidas en
la Tabla \ref{tab:tests}, y pruebas manuales de los cinco casos de uso contra la base
real, cuya evidencia se concentra en la Tabla \ref{tab:evidencia}.

\begin{table}[H]
\caption{Pruebas unitarias (Vitest): 26 pruebas, 2 archivos, 0 fallos.}
\label{tab:tests}
\centering
\renewcommand{\arraystretch}{1.05}
\begin{tabular}{@{}L{0.38\columnwidth}L{0.54\columnwidth}@{}}
\toprule
\textbf{Archivo} & \textbf{Cobertura} \\
\midrule
\nolinkurl{tests/sello-qr.test.ts} & Formato y unicidad del folio; la firma detecta cambios en folio, comunidad, pueblo, técnica o productor \\
\nolinkurl{tests/comisiones.test.ts} & Reparto 800 $\rightarrow$ 720/80; el envío no lleva comisión; no se pierden centavos; lote B2B de 30 piezas a \$1\,500 (total \$38\,250, anticipo \$11\,475, saldo \$26\,775) \\
\bottomrule
\end{tabular}
\end{table}

\begin{table}[H]
\caption{Verificación manual de los cinco casos de uso contra la base real.}
\label{tab:evidencia}
\centering
\renewcommand{\arraystretch}{1.08}
\begin{tabular}{@{}L{0.07\textwidth}L{0.26\textwidth}L{0.31\textwidth}L{0.28\textwidth}@{}}
\toprule
\textbf{CU} & \textbf{Pantalla} & \textbf{Artefacto generado} & \textbf{Evidencia en la base} \\
\midrule
CU-01 & \texttt{/registro-productor} y \texttt{/promotor} & Usuario, perfil y verificación & \nolinkurl{9da68be6-cce2-4a44-9a72-0c8509a75a80} creado \texttt{pendiente} y aprobado \\
CU-04 & \texttt{/productor/publicar} & Producto, ficha, folio y firma & \nolinkurl{1fb66f12-d382-4c05-ba3c-d56fd13e120f}, folio \nolinkurl{RMX-PUE-2026-1EA2E8} \\
CU-06 & \texttt{/verificar/[folio]} & Verificación pública del sello & \texttt{verificado: true} en \nolinkurl{RMX-PUE-2026-3F9A1C} y \nolinkurl{RMX-PUE-2026-1EA2E8} \\
CU-07 & \texttt{/catalogo} & Filtros y mapa de procedencia & Mapa: Puebla 4 productores / 7 piezas; Oaxaca 1 / 1 \\
CU-08 & \texttt{/comprador/carrito} & Pedido, pago, liquidación y bitácora & \texttt{PED-2026-C1C0A80D}; \$350 + \$35 = \$385 + \$150 envío = \$500 \\
\bottomrule
\end{tabular}
\end{table}

% =========================== VI. CONCLUSIONES ============================
\section{Conclusiones}

Quedaron entregados los dos productos de la primera entrega. El modelo
entidad-relación está implementado en el motor con 17 tablas, 10 enumeraciones, 2
vistas, una función transaccional y ocho disparadores, con correspondencia
verificable uno a uno contra las 24 clases del diagrama UML, y con las reglas que en
un diseño convencional residirían en la aplicación resueltas por la propia base:
obligatoriedad de los nueve campos de origen, exclusividad del productor aprobado
para publicar, conservación del reparto del dinero e inalterabilidad de la bitácora.

Sobre ese modelo opera la versión~1.0 del sistema, con los cinco casos de uso
funcionando contra la base real y no contra datos en memoria. La evidencia quedó
registrada con identificadores concretos ---folio de pedido, folio de sello y
productor identificado---, de modo que cada afirmación del documento es
reproducible por el equipo.

El avance más relevante para el propósito del proyecto es que el origen de cada
pieza dejó de ser una declaración y pasó a ser un dato verificable por cualquier
persona, sin cuenta y sin conocimiento previo de la plataforma. La siguiente
versión incorpora el inicio de sesión con Supabase Auth (la columna
\texttt{usuarios.auth\_id} ya está lista), el canal mayorista por lote —con sus
tablas y su cálculo ya resueltos— y la pasarela de pago real.

% ============================ REFERENCIAS ================================
\begin{thebibliography}{99}
\sloppy
\bibitem{propuesta} Equipo 10, ``RaicesMX: Propuesta inicial de proyecto de software,'' Facultad de Ciencias de la Computación, BUAP, oct. 2026.
\bibitem{uml} Equipo 10, ``Diagrama de clases UML de RaicesMX (24 clases),'' Facultad de Ciencias de la Computación, BUAP, 2026.
\bibitem{jornadaOriente} La Jornada de Oriente, ``Coyotes llegan a revender hasta el triple de precio los productos de artesanos o campesinos,'' 2 sep. 2021. [En línea]. Disponible en: \url{https://www.lajornadadeoriente.com.mx/puebla/coyotes-precio-los-productos-de-artesanos/}
\bibitem{inegiCultura} INEGI, ``Cuenta Satélite de la Cultura de México 2024,'' comunicado de prensa, 2025. [En línea]. Disponible en: \url{https://www.inegi.org.mx/contenidos/saladeprensa/boletines/2025/cultura/CSCM2024_CP.pdf}
\bibitem{stripe} Stripe, ``Tarifas y comisiones (México),'' 2026. [En línea]. Disponible en: \url{https://stripe.com/mx/pricing}
\bibitem{lfppc} Cámara de Diputados del H. Congreso de la Unión, ``Ley Federal de Protección del Patrimonio Cultural de los Pueblos y Comunidades Indígenas y Afromexicanas,'' DOF 17-01-2022. [En línea]. Disponible en: \url{https://www.diputados.gob.mx/LeyesBiblio/pdf/LFPPCPCIA.pdf}
\bibitem{rfc2104} H. Krawczyk, M. Bellare, y R. Canetti, ``HMAC: Keyed-Hashing for Message Authentication,'' RFC 2104, feb. 1997. [En línea]. Disponible en: \url{https://www.rfc-editor.org/rfc/rfc2104}
\bibitem{owasp} OWASP Foundation, \textit{OWASP Top 10:2021}, 2021. [En línea]. Disponible en: \url{https://owasp.org/Top10/}
\bibitem{rowlevel} Supabase, \textit{PostgreSQL Database Documentation: Row Level Security}, 2026. [En línea]. Disponible en: \url{https://supabase.com/docs/guides/database/postgres/row-level-security}
\end{thebibliography}

\end{document}
